\documentclass[10pt,a4paper]{amsart}
\usepackage[margin=19mm]{geometry}
\usepackage{amsmath,amssymb,mathtools}
\usepackage[scr=rsfs,cal=euler]{mathalfa}
\usepackage{xfrac}
\usepackage[unicode,hidelinks,bookmarks=false]{hyperref}
\newcommand{\ie}{i.e.\ }
\newcommand{\cf}{cf.\ }
\newcommand{\resp}{respectively\ }
\newcommand{\Subsection}{\subsection}
\providecommand{\nobreakdash}{\nobreak\hbox{-}}
\setlength{\emergencystretch}{2em}
\multlinegap0pt
\usepackage[shortlabels]{enumitem}
\usepackage{amssymb}
\usepackage{tikz}
\newtheorem{lem}{Lemma}
\newtheorem{thm}{Theorem}
\newcommand{\R}{\mathbb{R}}
\newcommand{\chara}[1]{\chi_{#1}}
\newcommand{\lowers}{LSC}
\newcommand{\mcal}{\mathcal{M}}
\newcommand{\tcal}{\mathcal{T}}
\title{Fully nonlinear oblique transmission problems: Well-posedness}
\author{I\~nigo U. Erneta}
\date{Source: arXiv:2606.14856v2 (CC BY 4.0)}
\begin{document}
\maketitle

\noindent\emph{Source and licence.} Fully nonlinear oblique transmission problems: Well-posedness. Version arXiv:2606.14856v2 (CC BY 4.0), distributed by arXiv under the Creative Commons Attribution 4.0 International licence, http://creativecommons.org/licenses/by/4.0/, as recorded in the arXiv licence metadata for this exact version and shown on the abstract page https://arxiv.org/abs/2606.14856v2. Full licence notice: \texttt{licenses/arxiv-2606.14856-CC-BY-4.0.txt}.\par\smallskip

\noindent\textit{Editorial scope.} Counted unit: the article's thm thm:abp (source0.tex lines 1092--1103). The article's complete original proof of it is delivered with it (source0.tex lines 1105--1236, one \texttt{proof} environment of 132 lines). No mathematics is written, repaired or re-derived: the statement and the proof are the article's own lines, in the article's own notation, and no retained line is substituted or re-flowed.  Retained uncounted as the source-local context the proof needs: lines 783--786; lines 862--875.  Nothing was omitted from any retained span, so the package carries no omission record.  No retained line contains a citation, so the wrapper carries no bibliography and no bibliography entry.  No retained line contains a figure environment, an image-inclusion command or drawing code, so no image is delivered and the package declares no asset dependency.  Editorial formatting only, in addition: an AMS \texttt{amsart} preamble that restates the article's own declarations and macros used by the retained lines, namely the environments \texttt{lem}, \texttt{thm} on separate counters, together with the standard packages those lines need.  The retained environments print in this document as Lemma 1, Theorem 1 in order of appearance, whereas the article prints the same items with the numbering of its own full document.  The complete native source of the article as distributed in the arXiv e-print for this exact version is bundled unchanged as \texttt{sources/202810/source0.tex}.
\begin{equation}
\label{tmin}
\tcal^{-}(\xi^{+}, \xi^{-}; \lambda, \Lambda, \mu) = \lambda \xi_n^{+} - \Lambda \xi_n^{-} - 2\mu |\xi'|,
\end{equation}

\begin{lem}
\label{lem:aux}
Let $u \in \underline{S}(\lambda,\Lambda,\mu,\theta; f^{+}, f^{-}, g)$ in $\Omega$, and let $\varphi$ be an admissible function satisfying
\[
\begin{cases}
\mcal^{+}(D^2 \varphi; \frac{\lambda}{n}, \Lambda) \leq \widetilde{f}^{\pm}(x) & \text{ in } \Omega^{\pm},\\
\tcal^{+}(\nabla \varphi^{+}, \theta \nabla \varphi^{-}; \lambda, \Lambda, \mu) \leq \widetilde{g}(x) & \text{ on } T.
\end{cases}
\]
Then
\[
u - \varphi \in \underline{S}(\lambda,\Lambda,\mu,\theta; f^{+}- \widetilde{f}^{+}, f^{-}- \widetilde{f}^{-}, g - \widetilde{g}) \quad \text{ in } \Omega.
\]
\end{lem}

\begin{thm}[ABP estimate]
\label{thm:abp}
Let $u \in \overline{S}(\lambda,\Lambda,\mu, \theta; f^{+}, f^{-}, g)$ in $B_R$, with $u \in \lowers(\overline{B_R})$.

There is a universal constant $C > 0$ such that
\[
\sup_{B_R} u_{-} \leq \sup_{\partial B_R} u_{-} + C R \left( 
 \|f^{+}_{+}\|_{L^{n}(\{u = \Gamma_u\} \cap B_R^{+})} + \|f^{-}_{+}\|_{L^{n}(\{u = \Gamma_u\} \cap B_R^{-})}+ \|g_{+}\|_{L^{\infty}(T_R)}
 \right),
\]
where $\Gamma_u$ denotes the convex envelope of $- u_{-} \chara{B_R}$ in the ball $B_{2R}$.
\end{thm}

\begin{proof}
Consider the piecewise linear barrier
\begin{equation}
\label{abp:bar}
\varphi(x) := \frac{1}{\lambda} \|g_{+}\|_{L^{\infty}(T_R)} (x_n)_{+}.
\end{equation}
Since $\varphi_{x_n}^{+} \geq 0$ and $\varphi_{x_n}^{-} = 0$, from~\eqref{tmin}, we have that
\[
\tcal^{-}(\nabla \varphi^{+},\theta \nabla \varphi^{-}; \lambda, \Lambda, \mu)
= \lambda \varphi_{x_n}^{+}
=  \|g_{+}\|_{L^{\infty}(T_R)},
\]
and hence by Lemma~\ref{lem:aux}
\begin{equation}
\label{abp:sup}
u- \varphi \in \overline{S}(\lambda, \Lambda, \mu, \theta; f^{+}, f^{-}, 0) \quad \text{ in } B_R.
\end{equation}

\medskip

For $\xi = (\xi', \xi_n) \in \R^{n-1} \times (0,\infty)$, let $\ell_{\xi}$ be the continuous piecewise linear function
\begin{equation}
\label{abp:piece}
\ell_{\xi} 
= \xi' \cdot x' + \xi_n\Big\{ 2\frac{\Lambda}{\lambda} (x_n)_{+} - (x_n)_{-}\Big\}.
\end{equation}
Since $(\ell_{\xi})_{x_n}^{+} = 2 \frac{\Lambda}{\lambda}\xi_n >0$ and $(\ell_{\xi})_{x_n}^{-} = \xi_n >0$, recalling $0 \leq \theta \leq 1$, from~\eqref{tmin} we have that
\[
\tcal^{-} (\nabla \ell_{\xi}^{+}, \theta \nabla \ell_{\xi}^{-}; \lambda, \Lambda, \mu) 
= \lambda(\ell_{\xi})_{x_n}^{+} - \Lambda (\ell_{\xi})_{x_n}^{-} - \mu(1+\theta) |\nabla' \ell_{\xi}|
\geq \Lambda \xi_n - 2 \mu |\xi'|.
\]
Therefore, if a vertical translation of $\ell_{\xi}$ touches $u-\varphi$ from below somewhere at $T_R$, then 
\begin{equation}
\label{abp:trans}
\Lambda \xi_n - 2\mu |\xi'| \leq 0
\end{equation}
by definition of viscosity supersolution.

\medskip

Let $v$ be the auxiliary function
\begin{equation}
\label{abp:v}
v := u- \varphi + \sup_{\partial B_R} (u-\varphi)_{-}
\end{equation}
and observe that $v \geq 0$ on $\partial B_R$.
We assume that its negative part is nontrivial, with
\begin{equation}
\label{abp:mdef}
M := \sup_{B_R} v_{-} > 0.
\end{equation}
Let $c_{\xi} \in \R$ be such that the vertical translation $\widetilde{\ell}_{\xi} := \ell_{\xi} + c_{\xi}$ touches $- v_{-} \chara{B_R}$ from below on $\overline{B}_{2R}$.
Finally, recall that $\Gamma_{v}$ is the convex envelope of $- v_{-} \chara{B_R}$ in $B_{2R}$, i.e., by definition
\[
\Gamma_v(x) = \sup\big\{
\ell(x) \colon \quad \ell \text{ affine}, \quad \ell \leq - v_{-} \chara{B_R} \, \text{ in } B_{2R}
\big\}.
\]

\medskip

We note the following properties of $\widetilde{\ell}_{\xi}$ depending on the dimensions of the parameter $\xi$:
\begin{enumerate}[label=(\roman*)]
\item If $|\xi| < \frac{\lambda}{6 \Lambda} \frac{M}{R}$, then $|\nabla \ell_{\xi}^{\pm}| < \frac{M}{3R}$ and hence $\widetilde{\ell}_{\xi}$ touches $-v_{-} \chara{B_R}$ from below in $B_R$.
\item If $\xi_n > \frac{2 \mu}{\Lambda} |\xi'|$, then by~\eqref{abp:trans} and the discussion above, the function $\widetilde{\ell}_{\xi}$ touches $-v_{-}\chara{B_R}$ from below in $\overline{B}_{2R} \setminus T_R$, i.e., away from the interface~$T_R$.
\end{enumerate}
Motivated by conditions (i) and (ii), we define the coefficient set 
\[
\Xi := \left\{\xi = (\xi', \xi_n) \in \R^{n-1}\times \R \colon \quad  \xi_n > \frac{2\mu}{\Lambda} |\xi'|, \quad |\xi| < \frac{\lambda}{6 \Lambda}\right\}.
\]
By construction, each $\xi \in \frac{M}{R}\Xi$ yields a piecewise linear function $\widetilde{\ell}_{\xi}$ below $- v_{-} \chara{B_R}$ in $\overline{B}_{2R}$ and touching $-v_{-}$ in $B_R\setminus T_R$.
Therefore, recalling definition~\eqref{abp:piece}, we have the following alternative depending on whether such a tangency point lies in $B_R^{-}$ or in $B_R^{+}$:
\[
\nabla \widetilde{\ell}^{-}_{\xi} = \xi 
\in \nabla \Gamma_{v}(B_R \setminus T_R)
\qquad \text{ or } \qquad 
\nabla \widetilde{\ell}^{+}_{\xi} = A (\xi) \in \nabla \Gamma_{v}(B_R \setminus T_R),
\]
where $A \colon \R^{n} \to \R^{n}$ is the linear map given by
\[
A(\xi', \xi_n) = (\xi', \textstyle\frac{2\Lambda}{\lambda} \xi_n) \qquad \text{ for } (\xi', \xi_n) \in \R^{n-1} \times \R.
\]

\medskip

We partition the rescaled coefficients by $\frac{M}{R}\Xi = X_1 \cup X_2$, where
\[
X_1 := \textstyle\frac{M}{R}\Xi
\cap \nabla \Gamma(B_R \setminus T_R)
\qquad \text{ and } \qquad 
X_2 := \textstyle \frac{M}{R}\Xi \setminus X_1.
\]
From the dichotomy above, it follows that
\[
X_1 \cup A(X_2) \subset \nabla \Gamma_{v}(B_R \setminus T_R)
\]
and taking the Lebesgue measure on both sides, we obtain the lower bound
\begin{equation}
\label{eq:trinity}
|\nabla \Gamma_{v}(B_R \setminus T_R)| \geq \max\left\{|X_1|, |A(X_2)| \right\} \geq \max\left\{|X_1|, |X_2| \right\} \geq \frac{1}{2}\left|\frac{M}{R}\Xi\right| = \frac{1}{2} |\Xi| \left(\frac{M}{R}\right)^{n}.
\end{equation}
On the other hand, since $\Gamma_v$ is $C^{1,1}$ away from the interface and recalling~\eqref{abp:sup} and~\eqref{abp:v}, by the classical argument we deduce an upper bound
\begin{equation}
\label{eq:trin2}
|\nabla \Gamma_{v}(B_R \setminus T_R)|
\leq \int_{\{v = \Gamma_v\}\cap (B_R\setminus T_R)} D^2 \Gamma_{v}
\leq C \left( \int_{\{u = \Gamma_u\}\cap B_R^{+}} (f_{+}^{+})^{n} 
+ \int_{\{u = \Gamma_u\}\cap B_R^{-}} (f_{+}^{-})^{n}
\right),
\end{equation}
where in the last line we used that $\{v = \Gamma_v\} \subset \{u = \Gamma_u\}$, a fact which follows by convexity of $\varphi$.
Since $\Xi$ has universal positive measure,~\eqref{eq:trinity} and~\eqref{eq:trin2} yield
\begin{equation}
\label{abp:puerta}
M \leq C R \left(\|f_{+}^{+}\|_{L^{n}(\{u = \Gamma_{u}\}\cap B_R^{+} )} + \|f_{+}^{-}\|_{L^{n}(\{u = \Gamma_{u}\}\cap B_R^{-} )}\right).
\end{equation}

\medskip

Combining~\eqref{abp:bar},~\eqref{abp:v},~\eqref{abp:mdef}, and~\eqref{abp:puerta}, we finally deduce
\[
\begin{split}
\sup_{B_R} u_{-}
&\leq \sup_{\partial B_R} (u-\varphi)_{-} + \sup_{B_R} v_{-}\\
&\leq \sup_{\partial B_R} u_{-}+ C R \|g_{+}\|_{L^{\infty}(T_R)} + C R \left(  
\|f_{+}^{+}\|_{L^{n}(\{u = \Gamma_{u}\}\cap B_R^{+} )} + \|f_{+}^{-}\|_{L^{n}(\{u = \Gamma_{u}\}\cap B_R^{-} )}
\right),
\end{split}
\]
which was the claim.
\end{proof}

\end{document}
